\documentclass[11pt]{article}
\usepackage[margin=1in]{geometry}
\usepackage{amsmath,amssymb}
\usepackage{booktabs}
\usepackage{longtable}

\title{Hardware-First Starlink Pass Simulation: Formula Report}
\author{Constellation Take-Home}
\date{April 2026}

\newcommand{\vect}[1]{\mathbf{#1}}
\newcommand{\norm}[1]{\left\lVert #1 \right\rVert}
\newcommand{\clip}{\operatorname{clip}}
\newcommand{\argmax}{\operatorname*{arg\,max}}
\newcommand{\sgn}{\operatorname{sgn}}
\newcommand{\atanTwo}{\operatorname{atan2}}

\begin{document}
\maketitle

\section{Purpose}

This document defines the simulation mathematically before any new implementation work.
The scope is a single bidirectional pass of one real Starlink satellite over one fixed
ground station. The implementation will precompute a discrete sequence of frames, and
every future visualization will read only from those frames.

The orbital propagator is the only external closed-form block:
\begin{equation}
  \left( \vect{r}^{ECI}_k, \vect{v}^{ECI}_k \right)
  =
  \operatorname{SGP4}(\mathrm{TLE}, t_k).
\end{equation}
This report treats the SGP4 outputs as the upstream state inputs and makes every
downstream equation explicit.

\section{Simulation Sequence}

The simulation runs at a fixed step size
\begin{equation}
  \Delta t = 1 \ \mathrm{s},
\end{equation}
over frame indices
\begin{equation}
  k \in \{0,1,\dots,K-1\},
  \qquad
  t_k = t_0 + k \Delta t.
\end{equation}

At each second, the per-frame dependency order is:
\begin{enumerate}
  \item Propagate orbital state with SGP4.
  \item Transform the satellite state into ECEF.
  \item Derive ground geometry: azimuth, elevation, range, range rate, Doppler.
  \item Evaluate operational visibility and the pass window.
  \item Evaluate active faults for the current second.
  \item Compute steering, phased-array degradation, RF losses, PA operating point,
        oscillator error, and ground-station transmit/receive pointing.
  \item Compute downlink and uplink link budgets, effective SNR, EVM, MCS, BER, and PER.
  \item Generate deterministic sampled packet traces for both directions.
  \item Aggregate retransmissions, latency, jitter, loss, and goodput for both directions.
  \item Update dynamic states for the next second: temperatures, battery energy,
        compute queue, and traffic queues.
\end{enumerate}

This ordering avoids circular dependencies: frame-$k$ outputs depend on the dynamic state
at the start of second $k$, and the dynamic state is then advanced to $k+1$.

\section{Independent Parameters}

\subsection{Configuration Inputs}

\begin{longtable}{p{0.18\textwidth} p{0.22\textwidth} p{0.50\textwidth}}
\toprule
Symbol & Meaning & Source / current scenario \\
\midrule
\endhead
$\mathrm{TLE}$ & Two-line element set & Real Starlink TLE, currently STARLINK-1008 / NORAD 44714 \\
$t_0, K, \Delta t$ & Simulation clock & Current repo window: 2026-04-12 04:15:05 UTC to 2026-04-12 04:40:04 UTC at 1 s cadence \\
$\phi_g, \lambda_g, h_g$ & Ground-station latitude, longitude, altitude & Current repo: $25.9897376^\circ$, $-97.1847795^\circ$, 10 m \\
$E_{\min}$ & Minimum operational elevation & Current repo: $10^\circ$ \\
$M(\alpha)$ & Horizon mask vs azimuth & Piecewise-linear interpolation of configured mask samples \\
$\alpha_{ant,0}, E_{ant,0}$ & Initial ground antenna boresight azimuth/elevation & Current repo: $92^\circ$, $31^\circ$ \\
$G^{gs,rx}_{nom}, G^{gs,tx}_{nom}$ & Ground-station nominal boresight receive / transmit gains & Current repo has 43.5 dBi receive gain; uplink transmit gain defaults to the same aperture gain \\
$f_c^{dl}, B^{dl}$ & Downlink carrier frequency and bandwidth & Current repo: 18.2 GHz, 250 MHz \\
$f_c^{ul}, B^{ul}$ & Uplink carrier frequency and bandwidth & Current repo: 28.3 GHz, 250 MHz \\
$L^{dl}_{off}, L^{ul}_{off}$ & Offered load for downlink and uplink & Current repo has one shared offered load; bidirectional implementation should split this into per-direction loads, defaulting $L^{ul}_{off}=L^{dl}_{off}$ if omitted \\
$L^{dl}_{pkt}, L^{ul}_{pkt}$ & Packet size for downlink and uplink & Current repo has one shared packet size; bidirectional implementation should default both directions to that value \\
$N^{dl}_{retx,max}, N^{ul}_{retx,max}$ & Maximum retransmissions per direction & Current repo has one shared retransmission limit; bidirectional implementation should default both directions to that value \\
$\tau^{dl}_{retx}, \tau^{ul}_{retx}$ & Retransmission wait per direction & Current repo has one shared retransmission delay; bidirectional implementation should default both directions to that value \\
$s$ & Deterministic simulation seed & Current repo: 44714 \\
$\mathcal{F}$ & Scenario fault list & Time-bounded deterministic faults \\
$N_{elem,0}$ & Nominal active array elements & Current repo default: 256 \\
$P_{out,\min}, P_{out,nom}, P_{out,\max}$ & PA operating limits & Current repo defaults: -5 dBW, 6 dBW, 9 dBW \\
$IBO_{nom}$ & Nominal PA output backoff & Current repo default: 2 dB \\
$NF_0$ & Nominal receiver noise figure & Current repo default: 1.8 dB \\
\bottomrule
\end{longtable}

\subsection{Implementation Constants to Be Chosen Explicitly}

The following are independent model coefficients that the implementation must define in
one formula-focused module, not inline in UI code:

\begin{longtable}{p{0.20\textwidth} p{0.68\textwidth}}
\toprule
Symbol group & Meaning \\
\midrule
\endhead
$G_{elem}, q_{scan}, \Delta \phi_{LSB}$ & Element gain, scan-loss exponent, phase-quantization step \\
$\theta_{FOR}$ & Maximum electronic field of regard \\
$k_{NF}, L_{filter}, L_{conv}, L_{impl}$ & RF front-end thermal sensitivity and fixed losses \\
$k_{comp}, P_{1dB}$ & PA compression coefficient and 1 dB compression reference \\
$\alpha_f, \dot{f}_{drift}, \sigma_{t,0}, k_t$ & Oscillator thermal drift, frequency drift rate, jitter coefficients \\
$C_i, R_i, A_i, \eta_i$ & Thermal capacitance, thermal resistance, exposed area, and coupling coefficients for subsystem $i$ \\
$A_{sol}, \eta_{sol}, E_{batt,\max}, P_{batt,\max}$ & Solar and battery parameters \\
$C_{comp}, c_{beam}, c_{dop}, c^{dl}_{modem}, c^{ul}_{modem}, c_{tm}$ & Compute capacity and workload coefficients \\
$\dot{\alpha}_{max}, \dot{E}_{max}, \theta^{3dB}_{gs}$ & Ground-station azimuth slew rate, elevation slew rate, and shared aperture beamwidth used for both ground TX and RX pointing loss \\
$P^{gs}_{tx,nom}, P^{gs}_{tx,\min}, P^{gs}_{tx,\max}, L^{gs}_{tx,rf}, NF^{gs}_{nom}$ & Ground-station nominal uplink power, power limits, fixed TX RF loss, and nominal RX noise figure \\
$\sigma^{gs}_{P}, \sigma^{gs}_{G}, \sigma^{gs}_{NF}, \sigma^{gs}_{f}, \sigma^{gs}_{EVM}$ & Ground-station stochastic standard deviations for TX power, aperture gain, RX noise figure, reference offset, and residual transmitter impairment \\
$OH^{dl}_{mac}, OH^{ul}_{mac}, N^{dl}_{sample,max}, N^{ul}_{sample,max}, \tau_{proc}$ & Direction-specific MAC overhead, sample-trace caps, and fixed processing latency \\
$\{(\eta^{dl}_m, \gamma^{dl}_{req,m}, a^{dl}_m, BER^{dl}_{min,m}, BER^{dl}_{max,m})\}, \{(\eta^{ul}_m, \gamma^{ul}_{req,m}, a^{ul}_m, BER^{ul}_{min,m}, BER^{ul}_{max,m})\}$ & Direction-specific MCS tables for downlink and uplink \\
\bottomrule
\end{longtable}

\section{Dynamic State Variables}

The simulation carries the following dynamic state from frame to frame:
\begin{align}
  T^{arr}_k, \ T^{pa}_k, \ T^{rf}_k, \ T^{osc}_k
  & \quad \text{array, PA, RF, and oscillator temperatures}, \\
  \alpha^{gs}_k, \ E^{gs}_k
  & \quad \text{actual ground-station tracked boresight azimuth and elevation}, \\
  E^{batt}_k
  & \quad \text{battery energy}, \\
  q^{comp}_k
  & \quad \text{compute backlog in work units}, \\
  q^{dl}_k, \ q^{ul}_k
  & \quad \text{downlink and uplink network backlogs in packets}.
\end{align}

All other values in frame $k$ are algebraic outputs of the current dynamic state,
current propagated orbit state, and current configuration.

\section{Orbital State and Coordinate Transforms}

\subsection{Satellite State}

\begin{equation}
  \left( \vect{r}^{ECI}_k, \vect{v}^{ECI}_k \right)
  =
  \operatorname{SGP4}(\mathrm{TLE}, t_k).
\end{equation}

Let $\theta_{g,k}$ denote Greenwich sidereal angle at time $t_k$.
The standard ECI-to-ECEF transform is
\begin{align}
  \vect{r}^{ECEF}_k
  &= R_3(\theta_{g,k}) \vect{r}^{ECI}_k, \\
  \vect{v}^{ECEF}_k
  &= R_3(\theta_{g,k}) \vect{v}^{ECI}_k
     - \vect{\omega}_{\oplus} \times \vect{r}^{ECEF}_k,
\end{align}
where
\begin{equation}
  \vect{\omega}_{\oplus}
  =
  \begin{bmatrix}
    0 \\ 0 \\ \omega_{\oplus}
  \end{bmatrix},
  \qquad
  \omega_{\oplus} = 7.2921150 \times 10^{-5} \ \mathrm{rad/s}.
\end{equation}

\subsection{Ground-Station ECEF Position}

Using WGS84,
\begin{align}
  e^2 &= f(2-f), \\
  N_\phi &= \frac{a}{\sqrt{1 - e^2 \sin^2 \phi_g}},
\end{align}
and
\begin{equation}
  \vect{r}^{gs}
  =
  \begin{bmatrix}
    (N_\phi + h_g)\cos \phi_g \cos \lambda_g \\
    (N_\phi + h_g)\cos \phi_g \sin \lambda_g \\
    (N_\phi(1-e^2) + h_g)\sin \phi_g
  \end{bmatrix}.
\end{equation}

\section{Geometry, Visibility, and Pass Window}

\subsection{Relative Geometry}

The line-of-sight vector from ground station to satellite is
\begin{equation}
  \vect{\rho}_k = \vect{r}^{ECEF}_k - \vect{r}^{gs}.
\end{equation}

Project $\vect{\rho}_k$ into the local east-north-up frame:
\begin{equation}
  \begin{bmatrix}
    e_k \\ n_k \\ u_k
  \end{bmatrix}
  =
  \begin{bmatrix}
    -\sin \lambda_g & \cos \lambda_g & 0 \\
    -\sin \phi_g \cos \lambda_g & -\sin \phi_g \sin \lambda_g & \cos \phi_g \\
    \cos \phi_g \cos \lambda_g & \cos \phi_g \sin \lambda_g & \sin \phi_g
  \end{bmatrix}
  \vect{\rho}_k.
\end{equation}

Then
\begin{align}
  R_k &= \norm{\vect{\rho}_k}, \\
  \alpha_k &= \operatorname{mod}\left(\atanTwo(e_k, n_k), 2\pi \right), \\
  E_k &= \atanTwo\left(u_k, \sqrt{e_k^2+n_k^2} \right).
\end{align}

\subsection{Horizon Mask and Operational Visibility}

Let the configured horizon mask consist of azimuth-elevation samples
$\{(\alpha_i, E_i^{mask})\}$ sorted by azimuth.
For $\alpha_k$ between adjacent samples $i$ and $i+1$, define the piecewise-linear mask
\begin{equation}
  M(\alpha_k)
  =
  E_i^{mask}
  +
  \frac{\alpha_k - \alpha_i}{\alpha_{i+1} - \alpha_i}
  \left(E_{i+1}^{mask} - E_i^{mask}\right).
\end{equation}

The required elevation is
\begin{equation}
  E^{req}_k = \max(E_{\min}, M(\alpha_k)).
\end{equation}

Operational visibility is
\begin{equation}
  V_k = \mathbf{1}\{E_k \ge E^{req}_k\}.
\end{equation}

\subsection{Range Rate and Doppler}

Because the ground station is fixed in ECEF,
\begin{align}
  \dot{R}_k
  &= \frac{\vect{\rho}_k^\top \vect{v}^{ECEF}_k}{R_k}, \\
  f^{dop,dl}_k
  &= - \frac{\dot{R}_k}{c} f_c^{dl}, \\
  f^{dop,ul}_k
  &= - \frac{\dot{R}_k}{c} f_c^{ul}.
\end{align}

\subsection{Pass Window}

Define the explicit pass window
\begin{align}
  k_{in} &= \min\{k : V_k = 1\}, \\
  k_{out} &= \max\{k : V_k = 1\}, \\
  k_{apex} &= \argmax_{k:V_k=1} E_k.
\end{align}

The pass duration is
\begin{equation}
  T_{pass} = (k_{out} - k_{in} + 1)\Delta t.
\end{equation}

For the current repo state, the precomputed pass window inside the 1500-frame clock is:
\begin{align}
  k_{in} &= 518 \quad (2026\text{-}04\text{-}12\ 04{:}23{:}43 \ \mathrm{UTC}), \\
  k_{apex} &= 726 \quad (2026\text{-}04\text{-}12\ 04{:}27{:}11 \ \mathrm{UTC}), \\
  k_{out} &= 1001 \quad (2026\text{-}04\text{-}12\ 04{:}31{:}46 \ \mathrm{UTC}).
\end{align}

\section{Fault Activation}

Each deterministic fault $j \in \mathcal{F}$ is parameterized by a start time
$t^{start}_j$, a duration $\tau_j$, a type $\tau^{kind}_j$, and a severity $\sigma_j$.
Its activation state at second $k$ is
\begin{equation}
  F_{j,k}
  =
  \mathbf{1}\{t^{start}_j \le t_k < t^{start}_j + \tau_j\}.
\end{equation}

Any fault contribution to a parameter $x$ is applied additively or multiplicatively via
\begin{equation}
  \Delta x^{fault}_k = \sum_{j \in \mathcal{F}_x} F_{j,k} \sigma_j w_{x,j},
\end{equation}
where $w_{x,j}$ is the per-fault coupling coefficient for parameter $x$.

\section{Satellite Attitude, Steering, and Ground-Station Pointing}

\subsection{Satellite Steering Angle}

Assume a nadir-pointing spacecraft bus and an Earth-facing electronically steered array.
The array boresight unit vector is
\begin{equation}
  \hat{n}^{arr}_k = - \frac{\vect{r}^{ECEF}_k}{\norm{\vect{r}^{ECEF}_k}}.
\end{equation}

The line of sight from satellite to ground station is
\begin{equation}
  \hat{\ell}_k = \frac{\vect{r}^{gs} - \vect{r}^{ECEF}_k}{R_k}.
\end{equation}

The required steering angle is
\begin{equation}
  \theta^{steer}_k = \arccos \left( \hat{n}^{arr}_k \cdot \hat{\ell}_k \right).
\end{equation}

The field-of-regard gate is
\begin{equation}
  V^{FOR}_k = \mathbf{1}\{\theta^{steer}_k \le \theta_{FOR}\}.
\end{equation}

\subsection{Ground-Station Steering Mechanics}

The ground station mechanically tracks the satellite in azimuth and elevation.
The desired tracked boresight is taken directly from the current satellite look angles:
\begin{equation}
  \alpha^{cmd}_k = \alpha_k,
  \qquad
  E^{cmd}_k = E_k.
\end{equation}

Let the wrapped azimuth command error be
\begin{equation}
  \Delta \alpha^{cmd}_k
  =
  \operatorname{mod}\left(\alpha^{cmd}_k - \alpha^{gs}_k + \pi, 2\pi\right) - \pi,
\end{equation}
and the elevation command error be
\begin{equation}
  \Delta E^{cmd}_k = E^{cmd}_k - E^{gs}_k.
\end{equation}

With slew-rate limits $\dot{\alpha}_{max}$ and $\dot{E}_{max}$, the actual boresight
state evolves as
\begin{align}
  \alpha^{gs}_{k+1}
  &=
  \alpha^{gs}_k
  +
  \clip\left(
    \Delta \alpha^{cmd}_k,
    -\dot{\alpha}_{max}\Delta t,
    \dot{\alpha}_{max}\Delta t
  \right), \\
  E^{gs}_{k+1}
  &=
  E^{gs}_k
  +
  \clip\left(
    \Delta E^{cmd}_k,
    -\dot{E}_{max}\Delta t,
    \dot{E}_{max}\Delta t
  \right).
\end{align}

The initial tracking state is taken from configuration:
\begin{equation}
  \alpha^{gs}_0 = \alpha_{ant,0},
  \qquad
  E^{gs}_0 = E_{ant,0}.
\end{equation}

\subsection{Simple Stochastic Ground-Station Terminal Model}

The ground station is intentionally much simpler than the satellite, but it should still
be realistic. Rather than a detailed thermal/power/compute stack, it uses bounded
deterministic stochastic perturbations around nominal terminal behavior.

For each second, define a direction-independent pseudo-random source
\begin{align}
  x^{gs}_{k,0}
  &= \left(s + 2654435761k\right)\bmod 2^{32}, \\
  x^{gs}_{k,n+1}
  &= \left(1664525x^{gs}_{k,n} + 1013904223\right)\bmod 2^{32}, \\
  u^{gs}_{k,n}
  &= \frac{x^{gs}_{k,n}}{2^{32}},
\end{align}
and map it into bounded zero-mean perturbations
\begin{equation}
  \epsilon^{gs}_{k,n}
  =
  \clip\left(\sqrt{12}\left(u^{gs}_{k,n} - \tfrac{1}{2}\right), -2, 2\right).
\end{equation}

The ground-terminal per-second perturbations are then
\begin{align}
  \delta P^{gs}_k &= \sigma^{gs}_{P}\epsilon^{gs}_{k,1}, \\
  \delta G^{gs,rx}_k &= \sigma^{gs}_{G}\epsilon^{gs}_{k,2}, \\
  \delta G^{gs,tx}_k &= \sigma^{gs}_{G}\epsilon^{gs}_{k,3}, \\
  \delta NF^{gs}_k &= \sigma^{gs}_{NF}\epsilon^{gs}_{k,4}, \\
  \delta f^{gs}_k &= \sigma^{gs}_{f}\epsilon^{gs}_{k,5}, \\
  EVM^{2}_{gs,tx,k} &= \sigma^{gs}_{EVM}{}^{2}\epsilon^{gs}_{k,6}{}^{2}.
\end{align}

\subsection{Ground-Station Pointing Loss}

The antenna-pointing separation at the ground station is
\begin{equation}
  \theta^{gs}_k
  =
  \arccos \left(
    \sin E_k \sin E^{gs}_k
    +
    \cos E_k \cos E^{gs}_k \cos(\alpha_k - \alpha^{gs}_k)
  \right).
\end{equation}

For a shared aperture beamwidth $\theta^{3dB}_{gs}$, the pointing loss is
\begin{equation}
  L^{gs}_{point,k}
  =
  12 \left( \frac{\theta^{gs}_k}{\theta^{3dB}_{gs}} \right)^2.
\end{equation}

The effective ground-station gains are
\begin{equation}
  G^{gs,rx}_k = G^{gs,rx}_{nom} + \delta G^{gs,rx}_k - L^{gs}_{point,k},
  \qquad
  G^{gs,tx}_k = G^{gs,tx}_{nom} + \delta G^{gs,tx}_k - L^{gs}_{point,k}.
\end{equation}

\section{Phased Array, RF Chain, PA, and Oscillator}

\subsection{Active Elements and Array Gain}

Let $d^{elem}_k \in [0,1]$ be the degraded-element fraction after fault application.
Then
\begin{equation}
  N_{elem,k} = N_{elem,0}(1-d^{elem}_k).
\end{equation}

The ideal coherent array gain is
\begin{equation}
  G^{ideal}_{arr,k} = G_{elem} + 10 \log_{10}(N_{elem,k}).
\end{equation}

Model scan loss by
\begin{equation}
  L^{scan}_k = -10 \log_{10}\left( \cos^{q_{scan}} \theta^{steer}_k \right).
\end{equation}

If the phase shifter least-significant bit is $\Delta \phi_{LSB}$ and the random
phase-error standard deviation is $\sigma_{\phi,therm,k}$, then
\begin{align}
  \sigma^2_{\phi,k}
  &= \frac{\Delta \phi_{LSB}^2}{12} + \sigma_{\phi,therm,k}^2, \\
  L^{phase}_k
  &= -10 \log_{10}\left( e^{-\sigma^2_{\phi,k}} \right)
   = 4.343 \sigma^2_{\phi,k}.
\end{align}

Let $L^{arr}_{therm,k}$ denote thermal gain degradation and $L^{arr}_{point,k}$ denote
residual pointing loss. Then the effective array gain is
\begin{equation}
  G_{arr,k}
  =
  G^{ideal}_{arr,k}
  - L^{scan}_k
  - L^{phase}_k
  - L^{arr}_{therm,k}
  - L^{arr}_{point,k}.
\end{equation}

The same electronically steered aperture is used for satellite transmit and receive,
with optional calibration offsets:
\begin{equation}
  G^{sat,tx}_k = G_{arr,k} - \Delta G^{sat,tx},
  \qquad
  G^{sat,rx}_k = G_{arr,k} - \Delta G^{sat,rx}.
\end{equation}

\subsection{Beamwidth and Beam Quality}

Define an effective aperture width
\begin{equation}
  D^{eff}_{arr,k} = d_{elem} \sqrt{N_{elem,k}},
\end{equation}
where $d_{elem}$ is inter-element spacing in meters.

Then the approximate half-power beamwidth is
\begin{equation}
  \theta^{3dB}_{arr,k}
  = \kappa_{bw} \frac{\lambda}{D^{eff}_{arr,k}\cos \theta^{steer}_k},
  \qquad
  \lambda = \frac{c}{f_c}.
\end{equation}

A reviewable scalar beam-quality metric is
\begin{equation}
  Q^{beam}_k
  =
  \exp \left(
    -\sigma_{\phi,k}^2
    - \left(\frac{\theta^{arr}_{point,k}}{\theta^{3dB}_{arr,k}}\right)^2
  \right).
\end{equation}

\subsection{RF Front-End Loss and Noise Figure}

Let $T^{rf}_{ref}$ be the RF reference temperature.
The RF noise figure is
\begin{equation}
  NF_k = NF_0 + k_{NF}(T^{rf}_k - T^{rf}_{ref}) + \Delta NF^{fault}_k.
\end{equation}

For the bidirectional model, the satellite receive noise figure is
\begin{equation}
  NF^{sat,rx}_k = NF_k + \Delta NF^{sat,rx}.
\end{equation}

The total transmit-side RF loss is
\begin{equation}
  L^{tx}_{rf,k} = L_{filter} + L_{conv} + L_{impl} + \Delta L^{fault}_{rf,k}.
\end{equation}

\subsection{PA Command, Limits, Compression, and Efficiency}

Let $EIRP^{cmd}$ be the commanded target EIRP for the downlink. The requested PA output is
\begin{equation}
  P^{req}_{out,k}
  = EIRP^{cmd} - G_{arr,k} + L^{tx}_{rf,k}.
\end{equation}

Let the temperature- and bus-limited PA maximum be
\begin{equation}
  P^{lim}_{out,k}
  =
  P_{out,\max}
  - k^{pa}_T \max(0, T^{pa}_k - T^{pa}_{derate})
  - \Delta P^{bus}_k
  - \Delta P^{fault}_k.
\end{equation}

The applied PA output is
\begin{equation}
  P_{out,k}
  =
  \clip(P^{req}_{out,k}, P_{out,\min}, P^{lim}_{out,k}).
\end{equation}

With saturation reference $P_{sat}$, the input/output backoff is
\begin{equation}
  IBO_k = P_{sat} - P_{out,k}.
\end{equation}

The compression loss is modeled as
\begin{equation}
  L^{comp}_k
  =
  k_{comp}\max(0, IBO_{nom} - IBO_k)^2.
\end{equation}

The PA efficiency is
\begin{equation}
  \eta^{pa}_k
  =
  \clip\left(
    \eta^{pa}_0
    - k^{pa}_{\eta}(T^{pa}_k - T^{pa}_{ref})
    - \Delta \eta^{fault}_k,
    \eta^{pa}_{min},
    \eta^{pa}_{max}
  \right).
\end{equation}

Converting dBW to watts,
\begin{equation}
  P^{rf}_{W,k} = 10^{P_{out,k}/10},
\end{equation}
and the PA DC draw is
\begin{equation}
  P^{dc}_{pa,k} = \frac{P^{rf}_{W,k}}{\eta^{pa}_k}.
\end{equation}

\subsection{Oscillator Drift, Tracking Lag, and Residual Frequency Error}

Let $T^{osc}_{ref}$ be the oscillator reference temperature.
The oscillator absolute offset is
\begin{equation}
  \Delta f^{osc}_k
  =
  \Delta f^{osc}_0
  + \alpha_f (T^{osc}_k - T^{osc}_{ref})
  + \dot{f}_{drift}(t_k - t_0)
  + \Delta f^{fault}_k.
\end{equation}

Let the compute-induced tracking lag in frames be
\begin{equation}
  \ell_k = \left\lceil \frac{\tau^{sched}_k}{\Delta t} \right\rceil.
\end{equation}

The Doppler correction command uses a delayed estimate:
\begin{equation}
  \widehat{f}^{dop,dl}_k = f^{dop,dl}_{k-\ell_k},
  \qquad
  \widehat{f}^{dop,ul}_k = f^{dop,ul}_{k-\ell_k}.
\end{equation}

The residual carrier-frequency errors at the modem are
\begin{align}
  \Delta f^{res,dl}_k
  &=
  \Delta f^{osc}_k
  + f^{dop,dl}_k
  - \widehat{f}^{dop,dl}_k, \\
  \Delta f^{res,ul}_k
  &=
  \Delta f^{osc}_k
  + \Delta f^{gs,tx}
  + f^{dop,ul}_k
  - \widehat{f}^{dop,ul}_k.
\end{align}

The timing-jitter standard deviation is
\begin{equation}
  \sigma_{t,k}
  =
  \sigma_{t,0}
  + k_t (T^{osc}_k - T^{osc}_{ref})
  + \Delta \sigma^{fault}_{t,k}.
\end{equation}

\section{Thermal, Power, and Compute Dynamics}

\subsection{Thermal State Update}

For subsystem $i \in \{arr, pa, rf, osc\}$, let the thermal model be first-order:
\begin{equation}
  T^i_{k+1}
  =
  T^i_k
  +
  \frac{\Delta t}{C_i}
  \left(
    Q^i_{gen,k}
    + Q^i_{sun,k}
    - \frac{T^i_k - T_{sink}}{R_i}
  \right).
\end{equation}

The generated heat terms are
\begin{align}
  Q^{pa}_{gen,k}
  &= P^{dc}_{pa,k} - P^{rf}_{W,k}, \\
  Q^{arr}_{gen,k}
  &= P^{beam}_{0}V_k + k^{arr}_{scan}(\theta^{steer}_k)^2, \\
  Q^{rf}_{gen,k}
  &= P^{rf}_{0} + k^{rf}_{cpu}u^{cpu}_k, \\
  Q^{osc}_{gen,k}
  &= P^{osc}_{0}.
\end{align}

The sunlight heat term is
\begin{equation}
  Q^i_{sun,k} = \alpha_i A_i \Phi_{sun} S_k,
\end{equation}
where $S_k \in [0,1]$ is the sun-exposure factor.

\subsection{Solar, Battery, and Bus Budget}

Solar generation is
\begin{equation}
  P^{sol}_k = \eta_{sol} A_{sol} \Phi_{sun} S_k.
\end{equation}

Total electrical load is
\begin{equation}
  P^{load}_k
  =
  P^{dc}_{pa,k}
  + P^{comp}_k
  + P^{therm}_k
  + P^{house}_0.
\end{equation}

Battery energy evolves as
\begin{equation}
  E^{batt}_{k+1}
  =
  \clip\left(
    E^{batt}_k + \Delta t (P^{sol}_k - P^{load}_k),
    0,
    E_{batt,\max}
  \right).
\end{equation}

Battery state of charge is
\begin{equation}
  SOC_k = \frac{E^{batt}_k}{E_{batt,\max}}.
\end{equation}

The available bus power is
\begin{equation}
  P^{avail}_k
  =
  P^{sol}_k
  + P_{batt,\max}\mathbf{1}\{SOC_k > SOC_{min}\}.
\end{equation}

If load exceeds availability, the shedding factor is
\begin{equation}
  u^{shed}_k
  =
  \clip\left(
    \frac{P^{load}_k - P^{avail}_k}{P^{load}_k},
    0,
    1
  \right).
\end{equation}

The power-limit penalty applied to the PA command is
\begin{equation}
  \Delta P^{bus}_k = k^{bus}_P u^{shed}_k.
\end{equation}

\subsection{Compute Utilization and Scheduling}

Let the offered packet rates be
\begin{align}
  \lambda^{dl}_{off}
  &=
  \frac{L^{dl}_{off}}{8L^{dl}_{pkt}}, \\
  \lambda^{ul}_{off}
  &=
  \frac{L^{ul}_{off}}{8L^{ul}_{pkt}}.
\end{align}

The compute demand per second is
\begin{equation}
  D^{comp}_k
  =
  c_{idle}
  + c_{beam}N_{elem,k}V_k
  + c_{dop}V_k
  + c^{dl}_{modem}\lambda^{dl}_{off}L^{dl}_{pkt}
  + c^{ul}_{modem}\lambda^{ul}_{off}L^{ul}_{pkt}
  + c_{tm}|\{j : F_{j,k}=1\}|.
\end{equation}

The compute backlog update is
\begin{equation}
  q^{comp}_{k+1}
  =
  \max\left(0, q^{comp}_k + \Delta t(D^{comp}_k - C_{comp})\right).
\end{equation}

Utilization is
\begin{equation}
  u^{cpu}_k = \clip\left( \frac{D^{comp}_k}{C_{comp}}, 0, 1 \right).
\end{equation}

The schedule delay is
\begin{equation}
  \tau^{sched}_k = \frac{q^{comp}_k}{C_{comp}}.
\end{equation}

Compute electrical load is
\begin{equation}
  P^{comp}_k = P^{comp}_{idle} + (P^{comp}_{max} - P^{comp}_{idle})u^{cpu}_k.
\end{equation}

\section{Derived RF Output and Link Budget}

Unless marked with superscript $ul$, the RF and link equations in this section refer
to the downlink from satellite to ground station.

\subsection{Polarization and EIRP}

Let the downlink polarization mismatch loss be
\begin{equation}
  L^{pol,dl}_k
  =
  -10 \log_{10}\left(
    \left| \hat{e}^{sat,tx}_k{}^\dagger \hat{e}^{gs,rx} \right|^2
  \right).
\end{equation}

Then the downlink transmit EIRP is
\begin{equation}
  EIRP^{dl}_k
  =
  P_{out,k}
  + G^{sat,tx}_k
  - L^{tx}_{rf,k}
  - L^{comp}_k
  - L^{pol,dl}_k.
\end{equation}

\subsection{Downlink Path Loss, Noise, and Raw SNR}

The downlink free-space path loss is
\begin{equation}
  \lambda^{dl} = \frac{c}{f_c^{dl}},
  \qquad
  L^{fs,dl}_k = 20 \log_{10}\left( \frac{4\pi R_k}{\lambda^{dl}} \right).
\end{equation}

Model atmospheric attenuation with an elevation-dependent secant law:
\begin{equation}
  L^{atm,dl}_k
  =
  \frac{L^{zen,dl}}{\max(\sin E_k, \sin E_{floor})}.
\end{equation}

The downlink received carrier power is
\begin{equation}
  P^{rx,dl}_k
  =
  EIRP^{dl}_k
  + G^{gs,rx}_k
  - L^{fs,dl}_k
  - L^{atm,dl}_k
  - L^{misc,dl}.
\end{equation}

Ground receive equivalent noise temperature is
\begin{equation}
  NF^{gs,rx}_k = NF^{gs}_{nom} + \delta NF^{gs}_k,
  \qquad
  T^{eq,dl}_k = 290\left(10^{NF^{gs,rx}_k/10} - 1\right).
\end{equation}

Ground downlink system temperature is
\begin{equation}
  T^{sys,dl}_k = T^{ant,gs}_k + T^{eq,dl}_k.
\end{equation}

Downlink noise power over bandwidth $B^{dl}$ is
\begin{equation}
  N^{dl}_k = 10 \log_{10}\left( k_B T^{sys,dl}_k B^{dl} \right).
\end{equation}

The raw downlink SNR is
\begin{equation}
  \gamma^{raw,dl}_{k,dB} = P^{rx,dl}_k - N^{dl}_k,
  \qquad
  \gamma^{raw,dl}_{k,lin} = 10^{\gamma^{raw,dl}_{k,dB}/10}.
\end{equation}

\subsection{Downlink EVM and Effective SNR}

Thermal-noise EVM contribution:
\begin{equation}
  EVM^{2,dl}_{th,k} = \frac{1}{\gamma^{raw,dl}_{k,lin}}.
\end{equation}

Residual-frequency EVM contribution:
\begin{equation}
  EVM^{2,dl}_{cfo,k}
  =
  k^{dl}_{cfo}
  \left(\frac{\Delta f^{res,dl}_k}{R^{dl}_s}\right)^2,
\end{equation}
where $R^{dl}_s$ is the downlink symbol rate.

Timing-jitter contribution:
\begin{equation}
  EVM^{2,dl}_{jit,k} = \left( 2\pi B^{dl} \sigma_{t,k} \right)^2.
\end{equation}

PA-compression contribution:
\begin{equation}
  EVM^{2,dl}_{pa,k}
  =
  k^{dl}_{pa,evm}\max(0, P_{out,k} - P_{1dB})^2.
\end{equation}

Phase-error contribution:
\begin{equation}
  EVM^{2,dl}_{phase,k}
  =
  2\left(1 - e^{-\sigma_{\phi,k}^2}\right).
\end{equation}

The total downlink EVM is
\begin{equation}
  EVM^{2,dl}_k
  =
  EVM^{2,dl}_{th,k}
  + EVM^{2,dl}_{cfo,k}
  + EVM^{2,dl}_{jit,k}
  + EVM^{2,dl}_{pa,k}
  + EVM^{2,dl}_{phase,k}.
\end{equation}

Define the effective downlink SNR used by the modem as
\begin{equation}
  \gamma^{eff,dl}_{k,lin} = \frac{1}{EVM^{2,dl}_k},
  \qquad
  \gamma^{eff,dl}_{k,dB} = 10\log_{10}\gamma^{eff,dl}_{k,lin}.
\end{equation}

\subsection{Simple Uplink: Ground Station to Satellite}

For completeness, the uplink model uses a simple but realistic ground transmitter:
shared mechanical pointing, nominal terminal settings, and bounded stochastic variation
in TX power, gain, reference offset, and residual transmitter impairment. It does not
introduce a satellite-style ground thermal or scheduler stack.

The applied ground-station transmit power is
\begin{equation}
  P^{gs}_{tx,k}
  =
  \clip(P^{gs}_{tx,nom} + \delta P^{gs}_k, P^{gs}_{tx,\min}, P^{gs}_{tx,\max}).
\end{equation}

The uplink polarization mismatch loss is
\begin{equation}
  L^{pol,ul}_k
  =
  -10 \log_{10}\left(
    \left| \hat{e}^{gs,tx}{}^\dagger \hat{e}^{sat,rx}_k \right|^2
  \right).
\end{equation}

The uplink EIRP is
\begin{equation}
  EIRP^{ul}_k
  =
  P^{gs}_{tx,k}
  + G^{gs,tx}_k
  - L^{gs}_{tx,rf}
  - L^{pol,ul}_k.
\end{equation}

The uplink free-space and atmospheric losses are
\begin{equation}
  \lambda^{ul} = \frac{c}{f_c^{ul}},
  \qquad
  L^{fs,ul}_k = 20 \log_{10}\left( \frac{4\pi R_k}{\lambda^{ul}} \right),
\end{equation}
\begin{equation}
  L^{atm,ul}_k
  =
  \frac{L^{zen,ul}}{\max(\sin E_k, \sin E_{floor})}.
\end{equation}

The satellite uplink receive power is
\begin{equation}
  P^{rx,ul}_k
  =
  EIRP^{ul}_k
  + G^{sat,rx}_k
  - L^{fs,ul}_k
  - L^{atm,ul}_k
  - L^{misc,ul}.
\end{equation}

Satellite uplink receive noise is
\begin{equation}
  T^{eq,ul}_k = 290\left(10^{NF^{sat,rx}_k/10} - 1\right),
  \qquad
  T^{sys,ul}_k = T^{ant,sat}_k + T^{eq,ul}_k,
\end{equation}
\begin{equation}
  N^{ul}_k = 10 \log_{10}\left(k_B T^{sys,ul}_k B^{ul}\right).
\end{equation}

The raw uplink SNR is
\begin{equation}
  \gamma^{raw,ul}_{k,dB} = P^{rx,ul}_k - N^{ul}_k,
  \qquad
  \gamma^{raw,ul}_{k,lin} = 10^{\gamma^{raw,ul}_{k,dB}/10}.
\end{equation}

The uplink EVM budget is
\begin{align}
  EVM^{2,ul}_{th,k}
  &= \frac{1}{\gamma^{raw,ul}_{k,lin}}, \\
  EVM^{2,ul}_{cfo,k}
  &= k^{ul}_{cfo}\left(\frac{\Delta f^{res,ul}_k + \delta f^{gs}_k}{R^{ul}_s}\right)^2, \\
  EVM^{2,ul}_{jit,k}
  &= \left(2\pi B^{ul}\sigma_{t,k}\right)^2, \\
  EVM^{2,ul}_{phase,k}
  &= 2\left(1 - e^{-\sigma_{\phi,k}^2}\right).
\end{align}

With a simple stochastic ground-transmitter impairment term, the total uplink EVM is
\begin{equation}
  EVM^{2,ul}_k
  =
  EVM^{2,ul}_{th,k}
  + EVM^{2,ul}_{cfo,k}
  + EVM^{2,ul}_{jit,k}
  + EVM^{2,ul}_{phase,k}
  + EVM^{2}_{gs,tx,k},
\end{equation}
and the uplink effective SNR is
\begin{equation}
  \gamma^{eff,ul}_{k,lin} = \frac{1}{EVM^{2,ul}_k},
  \qquad
  \gamma^{eff,ul}_{k,dB} = 10\log_{10}\gamma^{eff,ul}_{k,lin}.
\end{equation}

\section{Adaptive Modulation, BER, and PER}

The downlink MCS table contains entries indexed by $m$, each with spectral efficiency
$\eta^{dl}_m$ and required SNR $\gamma^{dl}_{req,m}$ in dB.
The selected downlink MCS is
\begin{equation}
  m^{dl}_k
  =
  \max\{m : \gamma^{eff,dl}_{k,dB} \ge \gamma^{dl}_{req,m} + \Gamma^{dl}_{margin}\}.
\end{equation}

The downlink coded BER model is an explicit logistic curve around the selected MCS threshold:
\begin{equation}
  BER^{dl}_k
  =
  BER^{dl}_{min,m^{dl}_k}
  +
  \frac{BER^{dl}_{max,m^{dl}_k} - BER^{dl}_{min,m^{dl}_k}}
       {1 + \exp\left(a^{dl}_{m^{dl}_k}\left(\gamma^{eff,dl}_{k,dB} - \gamma^{dl}_{req,m^{dl}_k}\right)\right)}.
\end{equation}

For downlink packet length in bits
\begin{equation}
  L^{dl}_{bit} = 8L^{dl}_{pkt},
\end{equation}
the packet-error rate is
\begin{equation}
  PER^{dl}_k = 1 - (1 - BER^{dl}_k)^{L^{dl}_{bit}}.
\end{equation}

The scheduled downlink PHY payload service rate is
\begin{equation}
  C^{sched,dl}_k = \eta^{dl}_{m^{dl}_k} B^{dl} (1 - OH^{dl}_{mac}).
\end{equation}

The downlink service rate in packets per second is
\begin{equation}
  \mu^{dl}_k = \frac{C^{sched,dl}_k}{L^{dl}_{bit}}.
\end{equation}

The uplink modulation and coding selection is computed independently:
\begin{equation}
  m^{ul}_k
  =
  \max\{m : \gamma^{eff,ul}_{k,dB} \ge \gamma^{ul}_{req,m} + \Gamma^{ul}_{margin}\}.
\end{equation}

\begin{equation}
  BER^{ul}_k
  =
  BER^{ul}_{min,m^{ul}_k}
  +
  \frac{BER^{ul}_{max,m^{ul}_k} - BER^{ul}_{min,m^{ul}_k}}
       {1 + \exp\left(a^{ul}_{m^{ul}_k}\left(\gamma^{eff,ul}_{k,dB} - \gamma^{ul}_{req,m^{ul}_k}\right)\right)}.
\end{equation}

\begin{equation}
  L^{ul}_{bit} = 8L^{ul}_{pkt},
  \qquad
  PER^{ul}_k = 1 - (1 - BER^{ul}_k)^{L^{ul}_{bit}}.
\end{equation}

\begin{equation}
  C^{sched,ul}_k = \eta^{ul}_{m^{ul}_k} B^{ul} (1 - OH^{ul}_{mac}),
  \qquad
  \mu^{ul}_k = \frac{C^{sched,ul}_k}{L^{ul}_{bit}}.
\end{equation}

\section{Packet Trace, Latency, Jitter, Retransmissions, and Goodput}

The packet-trace model runs independently for each direction
$d \in \{dl, ul\}$. Below, $\lambda^{d}_{off}$, $\mu^d_k$, $q^d_k$, $PER^d_k$,
$L^{d}_{pkt}$, $\tau^{d}_{retx}$, and $N^{d}_{sample,max}$ are the direction-specific
quantities. The shared compute/scheduler delay $\tau^{sched}_k$ couples the two
directions through on-board resource contention.

\subsection{Network Backlog}

The traffic backlog update is
\begin{equation}
  q^{d}_{k+1}
  =
  \max\left(0, q^{d}_k + \Delta t(\lambda^{d}_{off} - \mu^{d}_k)\right).
\end{equation}

The queueing delay is
\begin{equation}
  \tau^{queue,d}_k
  =
  \begin{cases}
    q^{d}_k / \mu^{d}_k, & \mu^{d}_k > 0, \\
    +\infty, & \mu^{d}_k = 0.
  \end{cases}
\end{equation}

The propagation delay is
\begin{equation}
  \tau^{prop}_k = \frac{R_k}{c}.
\end{equation}

\subsection{Bounded Deterministic Packet Sampling}

The number of offered packets in one second is
\begin{equation}
  N^{offer,d}_k = \lambda^{d}_{off}\Delta t.
\end{equation}

The bounded sample count is
\begin{equation}
  N^{samp,d}_k = \min\left(N^{d}_{sample,max}, \lceil N^{offer,d}_k \rceil \right).
\end{equation}

Each sample packet represents
\begin{equation}
  w^{d}_k = \frac{N^{offer,d}_k}{N^{samp,d}_k}
\end{equation}
packets.

For deterministic sampling, define
\begin{align}
  x_{k,m,a,d,0}
  &= \left(
    s
    + 73856093k
    + 19349663m
    + 83492791a
    + \zeta_d
  \right)\bmod 2^{32}, \\
  x_{k,m,a,d,1}
  &= \left(1664525x_{k,m,a,d,0} + 1013904223\right)\bmod 2^{32}, \\
  u_{k,m,a,d}
  &= \frac{x_{k,m,a,d,1}}{2^{32}},
\end{align}
where $m$ is the sample-packet index and $a$ is the attempt index.

\subsection{Per-Packet Transmission Outcome}

The total number of allowed attempts is
\begin{equation}
  A^{d}_{max} = 1 + N^{d}_{retx,max}.
\end{equation}

For packet sample $(k,m)$, attempt $a$ fails if
\begin{equation}
  I^{fail,d}_{k,m,a} = \mathbf{1}\{u_{k,m,a,d} < PER^{d}_k\}.
\end{equation}

The first successful attempt index is
\begin{equation}
  a^{succ,d}_{k,m}
  =
  \min\{a \in \{1,\dots,A^{d}_{max}\} : I^{fail,d}_{k,m,a}=0\}.
\end{equation}

If no such attempt exists, the packet is dropped:
\begin{equation}
  D^{d}_{k,m} = \prod_{a=1}^{A^{d}_{max}} I^{fail,d}_{k,m,a}.
\end{equation}

The actual attempt count is
\begin{equation}
  n^{tx,d}_{k,m}
  =
  \begin{cases}
    a^{succ,d}_{k,m}, & D^{d}_{k,m}=0, \\
    A^{d}_{max}, & D^{d}_{k,m}=1.
  \end{cases}
\end{equation}

Successful packet latency is
\begin{equation}
  \tau^{d}_{k,m}
  =
  \tau^{prop}_k
  + \tau^{sched}_k
  + \tau^{queue,d}_k
  + \tau_{proc}
  + (n^{tx,d}_{k,m}-1)\tau^{d}_{retx}.
\end{equation}

\subsection{Aggregated Traffic Metrics}

Weighted packet-loss fraction:
\begin{equation}
  p^{loss,d}_k
  =
  \frac{\sum_{m=1}^{N^{samp,d}_k} w^{d}_k D^{d}_{k,m}}
       {\sum_{m=1}^{N^{samp,d}_k} w^{d}_k}.
\end{equation}

Weighted retransmission count:
\begin{equation}
  N^{retx,d}_k
  =
  \sum_{m=1}^{N^{samp,d}_k} w^{d}_k (n^{tx,d}_{k,m} - 1).
\end{equation}

Weighted goodput in bits/s:
\begin{equation}
  G^{d}_k
  =
  \frac{L^{d}_{bit}}{\Delta t}
  \sum_{m=1}^{N^{samp,d}_k} w^{d}_k (1-D^{d}_{k,m}).
\end{equation}

Weighted mean latency over successful packets:
\begin{equation}
  \bar{\tau}^{d}_k
  =
  \frac{
    \sum_{m=1}^{N^{samp,d}_k} w^{d}_k (1-D^{d}_{k,m})\tau^{d}_{k,m}
  }{
    \sum_{m=1}^{N^{samp,d}_k} w^{d}_k (1-D^{d}_{k,m})
  }.
\end{equation}

Weighted latency jitter:
\begin{equation}
  J^{d}_k
  =
  \sqrt{
    \frac{
      \sum_{m=1}^{N^{samp,d}_k} w^{d}_k (1-D^{d}_{k,m})(\tau^{d}_{k,m} - \bar{\tau}^{d}_k)^2
    }{
      \sum_{m=1}^{N^{samp,d}_k} w^{d}_k (1-D^{d}_{k,m})
    }
  }.
\end{equation}

\section{Explainability Tags}

Each frame should carry reason tags derived from normalized impairment scores:
\begin{align}
  S^{lowSNR,dl}_k &= \max(0, \gamma^{dl}_{req,m^{dl}_k} - \gamma^{eff,dl}_{k,dB}), \\
  S^{lowSNR,ul}_k &= \max(0, \gamma^{ul}_{req,m^{ul}_k} - \gamma^{eff,ul}_{k,dB}), \\
  S^{freq,dl}_k &= \frac{|\Delta f^{res,dl}_k|}{\Delta f^{dl}_{tol}}, \\
  S^{freq,ul}_k &= \frac{|\Delta f^{res,ul}_k|}{\Delta f^{ul}_{tol}}, \\
  S^{pa}_k &= \frac{L^{comp}_k}{L^{comp}_{tag}}, \\
  S^{array}_k &= \frac{L^{scan}_k + L^{phase}_k + L^{arr}_{therm,k}}{L^{array}_{tag}}, \\
  S^{thermal}_k &= \max_i \frac{T^i_k - T^i_{th}}{\Delta T^i_{tag}}, \\
  S^{compute}_k &= \frac{\tau^{sched}_k}{\tau^{deadline}}, \\
  S^{power}_k &= u^{shed}_k.
\end{align}

Tag $r$ is active when $S^r_k \ge 1$, and the dominant explanation is
\begin{equation}
  r^{dom}_k = \argmax_r S^r_k.
\end{equation}

\section{Per-Frame Outputs}

For review, each frame should contain at least the following dependent outputs:
\begin{enumerate}
  \item Geometry: $\alpha_k$, $E_k$, $R_k$, $\dot{R}_k$, $f^{dop,dl}_k$, $f^{dop,ul}_k$, $V_k$.
  \item Pass context: $(k_{in}, k_{apex}, k_{out})$ and in-pass / out-of-pass flags.
  \item Hardware: $\theta^{steer}_k$, $\alpha^{gs}_k$, $E^{gs}_k$, $\theta^{gs}_k$,
        $N_{elem,k}$, $G_{arr,k}$, $G^{sat,tx}_k$, $G^{sat,rx}_k$, $\theta^{3dB}_{arr,k}$,
        $Q^{beam}_k$, $NF_k$, $NF^{sat,rx}_k$, $P_{out,k}$, $IBO_k$, $L^{comp}_k$,
        $\Delta f^{osc}_k$, $\Delta f^{res,dl}_k$, $\Delta f^{res,ul}_k$, $\sigma_{t,k}$,
        temperatures, battery state, compute backlog.
  \item Link: $(EIRP^{dl}_k, P^{rx,dl}_k, N^{dl}_k, \gamma^{raw,dl}_k, EVM^{dl}_k, \gamma^{eff,dl}_k, m^{dl}_k, BER^{dl}_k, PER^{dl}_k)$ and $(EIRP^{ul}_k, P^{rx,ul}_k, N^{ul}_k, \gamma^{raw,ul}_k, EVM^{ul}_k, \gamma^{eff,ul}_k, m^{ul}_k, BER^{ul}_k, PER^{ul}_k)$.
  \item Traffic: sampled packet events plus $(p^{loss,dl}_k, N^{retx,dl}_k, \bar{\tau}^{dl}_k, J^{dl}_k, G^{dl}_k)$ and $(p^{loss,ul}_k, N^{retx,ul}_k, \bar{\tau}^{ul}_k, J^{ul}_k, G^{ul}_k)$.
  \item Explainability: active reason tags and dominant cause.
\end{enumerate}

\section{Implementation Organization}

To keep formulas reviewable, the implementation should mirror this document with pure
modules in one simulation layer:
\begin{enumerate}
  \item \texttt{simulation/geometry}: coordinate transforms, range, Doppler, visibility, pass window.
  \item \texttt{simulation/hardware}: steering, phased array, RF chain, PA, oscillator,
        thermal, power, and compute.
  \item \texttt{simulation/link}: direction-specific EIRP, path loss, noise, EVM,
        effective SNR, MCS, BER, and PER for both downlink and uplink.
  \item \texttt{simulation/traffic}: deterministic packet sampling, retransmissions,
        latency, jitter, and goodput for both directions.
  \item \texttt{simulation/build-frame}: ordered orchestration from propagated state to
        final frame outputs.
\end{enumerate}

No visualization code should contain simulation math. Every future UI surface should read
already-computed frame fields only.

\end{document}
